\documentclass[a4paper]{article}

\usepackage{graphicx}
\usepackage{amsmath,amssymb}
\usepackage{minted}
\usepackage{multirow}
\usepackage{float}
\usepackage{todonotes}
\usepackage{forest}
\usepackage{xstring}
\usepackage{xcolor}
\usepackage[most]{tcolorbox}
\usepackage{xparse}
\usepackage{natbib}

\usetikzlibrary{graphs,graphdrawing}
\usegdlibrary{layered}

% for external graphviz
\input{/home/donkarlo/Dropbox/repo/nd_ai_project/src/nd_ai/knoweldge_representation/language/formal/computer/programming/tex/latex/code_clips/external_graphviz.tex}

% for tags
\input{/home/donkarlo/Dropbox/repo/nd_ai_project/src/nd_ai/knoweldge_representation/language/formal/computer/programming/tex/latex/parts/control_sequence/kind/semtag/semtag.tex}

% for links
\input{/home/donkarlo/Dropbox/repo/nd_ai_project/src/nd_ai/knoweldge_representation/language/formal/computer/programming/tex/latex/code_clips/seehere.tex}

\title{Mind wandering}
\date{}

\begin{document}
\maketitle
    \seeherefooter{https://en.wikipedia.org/wiki/Default_mode_network}

    \cite{mason-2007-wandering-minds-the-default-network-and-stimulus-independent-thought} examine the relation between the default network and stimulus-independent thought, especially mind wandering. The paper argues that when people are not fully occupied by external tasks, internally generated thoughts become more likely. Using fMRI evidence, the authors link activity in the default network to these wandering, self-generated mental states. The paper is important because it connects the default mode network not only to rest, but also to spontaneous cognition and task-unrelated thought.
    \bibliographystyle{plainnat}
    \bibliography{/home/donkarlo/Dropbox/repo/nd_shared_research/bibliography/bibliography.bib}
\end{document}